\chapter{OCR and Translation Calibration Pilot Design}\label{ch:ocr-pilot}

This chapter specifies the pilot study design for calibrating OCR confidence
thresholds \emph{and} the Arabic$\to$English translation quality gate,
addressing the review findings that both the OCR thresholds (30\,\%/55\,\%)
and the Levenshtein distance gate ($\le$0.30) are calibrated on three
documents. Both calibrations are run against the same pilot corpus to avoid
duplicate ground-truth effort.

\section{Objective}

The pilot aims to:
\begin{enumerate}
  \item Determine empirically-validated OCR confidence thresholds per country
  \item Establish confidence intervals for threshold boundaries
  \item Identify document types that require special handling
  \item Create a baseline for ongoing threshold monitoring
  \item Calibrate the Arabic$\to$English translation quality gate
        (Levenshtein distance threshold) against the same pilot corpus
\end{enumerate}

\section{Pilot design}

\subsection{Sample size}

\begin{table}[htbp]
\centering
\caption{Pilot sample size by country.}
\label{tab:pilot-sample}
\footnotesize
\begin{tabularx}{\linewidth}{@{}l r r X@{}}
\toprule
\textbf{Country} & \textbf{Minimum} & \textbf{Target} & \textbf{Stratification} \\ \midrule
Saudi Arabia (SA) & 50 & 75 & By invoice type (tax, commercial) \\
Egypt (EG) & 50 & 75 & By scan quality (photo, scan, PDF) \\
Bahrain (BH) & 50 & 75 & By document language (AR, EN, mixed) \\
Qatar (QA) & 50 & 75 & By amount range (small, medium, large) \\
\textbf{Total} & \textbf{200} & \textbf{300} & \\
\bottomrule
\end{tabularx}
\end{table}

The sample size of 50 per country provides:
\begin{itemize}
  \item 95\% confidence interval width of $\pm$7\% for threshold estimates
  \item Statistical power of 80\% to detect 15\% differences between countries
  \item Sufficient stratification for 3--4 document types per country
\end{itemize}

\subsection{Document selection}

Documents are selected by:
\begin{enumerate}
  \item Random sampling from the last 6 months of invoices
  \item Stratified by document type (tax invoice, commercial invoice, payment
        instruction)
  \item Stratified by scan quality (high: $>$300 DPI, medium: 150--300 DPI,
        low: $<$150 DPI)
  \item Excluding documents with known data quality issues
\end{enumerate}

\subsection{Ground truth creation}

For each pilot document:
\begin{enumerate}
  \item Two independent human reviewers manually transcribe all fields
        in the source language (Arabic) and the required English fields
  \item Inter-rater agreement calculated (target: Cohen's $\kappa > 0.9$)
        for categorical fields; character error rate (CER) $\le$1\,\%
        on free-text fields
  \item Disagreements are resolved by a third reviewer whose decision is
        final. The resolution process is:
        \begin{enumerate}
          \item Third reviewer inspects both transcriptions side-by-side
                with the source document.
          \item Third reviewer selects one of the two transcriptions
                \emph{or} provides a new transcription; the choice and
                rationale are logged to
                \path{pilot/disputes/<doc-id>.json}.
          \item The third reviewer's verdict is the ground truth; the
                pilot lead reviews the top 10\,\% of disputes monthly
                to detect systematic bias.
        \end{enumerate}
  \item Final ground truth stored in structured JSON
        (\path{pilot/ground-truth/<country>/<doc-id>.json}), validated
        against \texttt{docs/schema/arabic/ocr-result.schema.json}
\end{enumerate}

\section{Threshold calibration methodology}

\subsection{OCR evaluation}

Each document is processed by the OCR pipeline:
\begin{enumerate}
  \item Record per-field OCR confidence scores
  \item Record per-document aggregate confidence
  \item Compare OCR output to ground truth
  \item Calculate field-level accuracy (character error rate)
\end{enumerate}

\subsection{Threshold determination}

For each country, determine optimal thresholds using:

\begin{lstlisting}[basicstyle=\ttfamily\small,caption={Threshold calibration algorithm.}]
def calibrate_thresholds(pilot_data, target_precision=0.95):
    """
    Calibrate OCR confidence thresholds.

    Parameters:
        pilot_data: List of (confidence, accuracy) pairs
        target_precision: Minimum acceptable precision (default 95%)

    Returns:
        high_threshold: Above this, auto-accept
        low_threshold: Below this, reject
    """
    # Sort by confidence descending
    sorted_data = sorted(pilot_data, key=lambda x: -x.confidence)

    # Find high_threshold: highest confidence where
    # precision >= target for all docs above threshold
    high_threshold = find_precision_cutoff(
        sorted_data, target_precision
    )

    # Find low_threshold: confidence below which
    # accuracy < 50% (not worth manual review)
    low_threshold = find_accuracy_floor(sorted_data, floor=0.50)

    return high_threshold, low_threshold
\end{lstlisting}

\paragraph{Justification of \texttt{target\_precision = 0.95}.}
The AP workflow routes any invoice below the high threshold to
two-person manual review (Chapter~5), which is tolerant of 5\,\%
false-positive auto-accepts from a \emph{business-risk} standpoint:
the downstream VAT engine (Chapter~4) independently validates
tax-ID and amount fields, and the reconciliation step catches
posting errors before payment is released. The 97\,\% soundness
threshold reported in the referenced OCR evaluation paper is
appropriate for \emph{unsupervised} pipelines; because our pipeline
has a manual review backstop, 95\,\% provides the right
precision/recall trade-off (empirically validated in the pilot
success criteria below). A sensitivity analysis at 93\,\%, 95\,\%,
and 97\,\% must be produced as part of the pilot report; the final
threshold is the largest value in \{0.93, 0.95, 0.97\} for which
the success criteria in Section~\ref{sec:pilot-success} hold.

\subsection{Translation quality gate calibration}
\label{sec:pilot-translation}

Using the same pilot corpus, the Arabic$\to$English translation
quality gate (currently Levenshtein distance $\le$0.30 against the
model's re-translation, calibrated on three documents) is
re-calibrated:

\begin{enumerate}
  \item For each ground-truthed pilot document, run the
        Arabic$\to$English$\to$Arabic round-trip translation and
        compute the normalized Levenshtein distance against the
        source Arabic.
  \item Compute the field-level translation accuracy against the
        human ground truth (exact match on numeric fields, BLEU-4
        $\ge$0.6 on free-text fields).
  \item Determine the Levenshtein threshold $\tau^{*}$ that
        maximises the F1 between ``round-trip gate passed'' and
        ``field-level translation accurate'' across the 200 pilot
        documents.
  \item Report the per-country threshold and the 95\,\% bootstrap
        CI ($\tau^{*}_{\mathrm{lower}}, \tau^{*}_{\mathrm{upper}}$).
        Deploy using the conservative upper bound.
\end{enumerate}

\noindent
The translation calibration adds no additional ground-truth effort
(the human transcriptions already contain the English fields) and
extends the pilot timeline by one week (added to
Section~\ref{sec:pilot-timeline}).

\subsection{Confidence interval calculation}

Bootstrap sampling (1,000 iterations) determines confidence intervals:
\begin{itemize}
  \item Report 95\% CI for high threshold
  \item Report 95\% CI for low threshold
  \item Use upper bound of CI for conservative deployment
\end{itemize}

\section{Success criteria}\label{sec:pilot-success}

The pilot succeeds when all of the following are true:
\begin{enumerate}
  \item $\geq$200 documents processed (50 per country minimum)
  \item Inter-rater agreement $\kappa > 0.9$ (categorical fields)
        and CER $\le$1\,\% (free-text fields)
  \item 95\,\% CI width $\leq \pm 5\,\%$ for all thresholds
  \item High threshold achieves $\geq$95\,\% precision
  \item Low threshold correctly rejects $\geq$90\,\% of unusable documents
  \item Translation gate F1 $\geq$0.85 at the selected Levenshtein
        threshold (see Section~\ref{sec:pilot-translation})
\end{enumerate}

\section{Fallback procedure}

If the pilot fails success criteria:
\begin{enumerate}
  \item Increase sample size by 50\% and repeat
  \item If still failing, segment by document type and calibrate separately
  \item If per-type calibration fails, default to manual review for all
        documents from that country until more data is collected
\end{enumerate}

\section{Timeline}\label{sec:pilot-timeline}

\begin{table}[htbp]
\centering
\caption{OCR and translation pilot timeline.}
\label{tab:pilot-timeline}
\footnotesize
\begin{tabularx}{\linewidth}{@{}l l l@{}}
\toprule
\textbf{Phase} & \textbf{Duration} & \textbf{Deliverable} \\ \midrule
Document selection & Week 1 & 300 documents identified \\
Ground truth creation & Weeks 2--5 & All documents transcribed (200 min) \\
OCR processing & Week 6 & Confidence scores recorded \\
Translation round-trip & Week 6 & Levenshtein distances recorded \\
Threshold calibration & Week 7 & Per-country OCR + translation thresholds \\
Validation & Week 8 & Held-out set validation \\
\textbf{Total} & \textbf{8 weeks} & Production-ready thresholds \\
\bottomrule
\end{tabularx}
\end{table}

\paragraph{Ground-truth effort justification.}
The original 6-week plan compressed ground-truth creation into
2~weeks (``all documents transcribed'') for 200 documents with
two reviewers --- 400 person-document reviews in 10 working days,
or 40 per reviewer per day. Empirically (internal benchmark
2025-Q4), a skilled Arabic-literate reviewer transcribes an
average invoice in 12--18~minutes, giving a realistic throughput
of 25--30 documents per 8-hour day with breaks and QA. The
revised 4-week ground-truth window (Weeks 2--5) supports
\emph{sustainable} throughput at 20~documents/reviewer/day, with
Week~5 reserved for dispute resolution (the third-reviewer loop
typically touches 8--12\,\% of documents). Pilot lead may
compress back to 3 weeks if a third reviewer is seconded for
dispute resolution only.

\section{Ownership}

\begin{table}[htbp]
\centering
\caption{Pilot roles and responsibilities.}
\label{tab:pilot-roles}
\footnotesize
\begin{tabularx}{\linewidth}{@{}l X@{}}
\toprule
\textbf{Role} & \textbf{Responsibility} \\ \midrule
Pilot Lead & Overall coordination, success criteria validation \\
Data Analyst & Sample selection, threshold calibration \\
QA Reviewers (2) & Ground truth creation, inter-rater agreement \\
ML Engineer & OCR pipeline configuration, metrics collection \\
Product Owner & Go/no-go decision, fallback approval \\
\bottomrule
\end{tabularx}
\end{table}

\section{Post-pilot monitoring}

After threshold deployment:
\begin{enumerate}
  \item Monthly accuracy review on random 50-document sample
  \item Alert if accuracy drops $>$5\% from pilot baseline
  \item Annual recalibration with new 200-document sample
  \item Per-country threshold adjustment if document mix changes
\end{enumerate}